\documentclass[11pt]{article}
\usepackage{fontspec}
\usepackage[french,english]{babel}
\usepackage{amsmath}
\usepackage{amssymb}
\usepackage{amsthm}
\usepackage{mathrsfs}
\usepackage{tikz-cd}
\usepackage{enumitem}

\newcommand{\printedpage}[1]{\marginpar{\footnotesize p.~#1}}
\newcommand{\msp}[1]{\marginpar{\footnotesize [#1]}}
\newcommand{\D}{\mathbb{D}}
\newcommand{\sfup}[1]{\textsf{\upshape #1}}
\newcommand{\Dia}{\sfup{Dia}}
\newcommand{\Hot}{\sfup{Hot}}
\newcommand{\HOT}{\sfup{HOT}}
\newcommand{\Ob}{\sfup{Ob}}
\newcommand{\Hom}{\sfup{Hom}}
\newcommand{\Cat}{$\mathcal{C}$\textit{at}}

\begin{document}

\printedpage{69}
% head: A. Grothendieck : Dérivateurs, Ch. I, Édition des Universités de Montpellier II et Paris VII

\textbf{Proposition 3.} \textit{Let}
\[
\begin{tikzcd}
X \arrow[rr, "f"] \arrow[dr, "p"'] & & X' \arrow[dl, "p'"] \\
& S &
\end{tikzcd}
\]
\textit{be in \Dia, with $p$ and $p'$ weak cohomological $\D$-fibrations} (resp. \textit{weak homological $\D$-fibrations}), \textit{let $S_0 \subset S$ be such that $\pi_0(S_0) \to \pi_0(S)$ is surjective. For $f$ to be a $\D$-equivalence, it suffices that the $f_s : X_s \to X'_s$, for $s \in S_0$, be so.}

\textbf{Proposition 4.} \textit{If $f : X \to Y$ in \Dia{} is $\D$-proper} (resp. $\D$-\textit{smooth}), \textit{then $f$ is a weak cohomological} (resp. \textit{homological}) \textit{$\D$-fibration if and only if for every $F$ constant on $Y$, $f_*f^*F$} (resp. $f_!f^*F$) \textit{is locally constant on $Y$. And if this is so, for every $Y' \to Y$ in \Dia, $f' : X' = X \times_Y Y' \to Y'$ is again a weak cohomological} (resp. \textit{homological}) \textit{$\D$-fibration.}

Immediate from criterion \textit{b}) of proposition 1.

\msp{page 117}\textbf{Proposition 5.} \textit{Let $f : X \to Y$ be in \Dia. Consider the following conditions:}

a) \textit{For every $Y'' \to Y'$ over $Y$ in \Dia, if $Y''$ and $Y'$ are \Hot-aspheric, then}
\[
X'' = X \times_Y Y'' \longrightarrow X' = X \times_Y Y'
\]
\textit{is a $\D$-equivalence.}
\[
\begin{tikzcd}
X & X' \arrow[l] & X'' \arrow[l] \\
Y \arrow[from=u] & Y' \arrow[l] \arrow[from=u] & Y'' \arrow[l, "p"] \arrow[from=u]
\end{tikzcd}
\]

a$'$) \textit{As in} a), \textit{but assuming $p : Y'' \to Y'$ is a \Hot-equivalence (instead of assuming $Y'$, $Y''$ \Hot-aspheric).}

a$_1$) \textit{As in} a), \textit{but assuming (instead of asphericity) that $Y'$ has a final object, and $Y''$ is the unit category.}

a$_2$) \textit{As in} a), \textit{but assuming that $Y'$ has a cofinal object, and $Y''$ is the unit category.}

b$_1$) \textit{$f$ is a weak cohomological $\D$-fibration, and remains so after any base change.}

b$_2$) \textit{$f$ is a weak homological $\D$-fibration, and remains so after any base change.}

\printedpage{70}
% head: A. Grothendieck : Dérivateurs, Ch. I, Édition des Universités de Montpellier II et Paris VII
\textit{We then have the diagram of tautological implications}
\[
\begin{tikzcd}[column sep=small]
& & & a' \arrow[d, Rightarrow] & & & & \\
& & & a \arrow[dlll, Rightarrow] \arrow[drrr, Rightarrow] & & & & \\
a_1 & \Longleftrightarrow & b_1 & & b_2 & \Longleftrightarrow & a_2 & .
\end{tikzcd}
\]
\textit{When $\D$ satisfies axiom} \textbf{Der} 3 a$'$ (p. 111), \textit{then all these conditions, except} a$'$), \textit{are equivalent. When $\D$ satisfies axiom} \textbf{Der} 3 b$'_1$ \textit{(or only} \textbf{Der} 3 b$''_{\mathrm{coh}}$ or b$''_{\mathrm{hom}}$\textit{),} \textit{then all the conditions (including} a$'$\textit{) are equivalent.}

To simplify, we shall assume henceforth that axiom \textbf{Der} 3 $a'$) is satisfied.

\textbf{Definition 3.} Under the equivalent conditions $a$, $a_1$, $b_1$, $a_2$, $b_2$ of proposition 4, we say that $f$ is a $\D$-\textit{fibration}. % [sic] proposition 4 -> proposition 5, the conditions a, a_1, b_1, a_2, b_2 are those of proposition 5

\msp{page 118}\textbf{Proposition 6.} \textit{Suppose that $f : X \to Y$ in \Dia{} is $\D$-proper} (resp. $\D$-\textit{smooth}). \textit{Equivalent conditions:}

a) \textit{$f$ is a $\D$-fibration.}

b) \textit{$f$ is a weak cohomological} (resp. \textit{homological}) \textit{$\D$-fibration.}

c) \textit{For every $y' \to y$ in $Y$, $X/y' \to X/y$} (resp. $y\backslash X \to y'\backslash X$) \textit{is a $\D$-equivalence.}

d) \textit{For every $F$ constant on $Y$, $f_*f^*F$} (resp. $f_!f^*F$) \textit{is locally constant on $Y$.}

d$'$) \textit{For every $F$ locally constant on $Y$, $f_*f^*F$} (resp. $f_!f^*F$) \textit{is locally constant on $Y$.}

This follows formally from Prop. 4 and Prop. 5.

\textbf{Corollary.} \textit{Suppose $f : X \to Y$ is} \Cat\textit{-fibrant} (resp. \Cat\textit{-cofibrant}). \textit{For $f$ to be a $\D$-fibration, it is necessary and sufficient that for every $u : y \to y'$ in $Y$, $u^* : X_{y'} \to X_y$} (resp. $u_* : X_y \to X_{y'}$) \textit{be a $\D$-equivalence.}

It is immediate that the notion of $\D$-fibration is stable under base change, but it is doubtful that it is so under composition.\footnote[60]{N.B. If the $\D$-equivalences are already \Hot-equivalences, then the notion coincides with that of \Hot-fibration, and it is stable under composition. Also, if the $\D$-equivalences are in the sense of cohomological, or homological, $\D$-equivalences, then the notion coincides with that of cohomological (resp. homological) $\D$-fibrations below, and is also stable under composition.} (Here again, I have omitted to make a counterexample, with derivators $\D = \sfup{Der}(k\text{-}\sfup{Mod})$, say.) Another shortcoming of the notion: it is not clear, if $f$ is a $\D$-fibration and $F$ a $\D$-coefficient on $Y$, whether
\[
f_*f^*F \qquad \text{and} \qquad f_!f^*F
\]
\printedpage{71}
% head: A. Grothendieck : Dérivateurs, Ch. I, Édition des Universités de Montpellier II et Paris VII
are computed ``fiber by fiber'', as one would expect in the case of a ``bundle''.

On the other hand, here is a nice result:

\msp{page 119}\textbf{Proposition 7.} \textit{Let $f : X \to Y$ be in \Dia. Equivalent conditions:}

a) \textit{$f$ is a ``universal'' $\D$-equivalence,} i.e. \textit{for every $Y' \to Y$ in \Dia, $f' : X' = X \times_Y Y' \to Y'$ is a $\D$-equivalence.}

b) \textit{$f$ is a $\D$-fibration with $\D$-aspheric fibers.}

\textbf{Definition 4.} Under the equivalent conditions of proposition 7, we say that $f$ is a \textit{trivial} $\D$-fibration.

But beware, $f$ can be a $\D$-fibration and a $\D$-equivalence, without the $\D$-fibration being trivial \textit{i.e.} without the fibers being $\D$-aspheric, cf. p. 87 and following. % [sic] peut-être -> peut être, verb "pouvoir" followed by infinitive
And also (should go above) even if $\D$ satisfies \textbf{Der} 3 $b_2$ (the strongest of the variants of \textbf{Der} 3), thus when \underline{all} the conditions of proposition 5, including $a'$), are equivalent, these (\textit{i.e.} that $f$ be a $\D$-fibration) nevertheless do not imply, in general, that in condition $a'$) one can restrict oneself to $p : Y'' \to Y'$ a $\D$-equivalence (instead of a \Hot-equivalence). In other words, if we have a cartesian diagram in \Dia
\[
\begin{tikzcd}
X \arrow[d, "f"'] & X' \arrow[l, "q"'] \arrow[d, "f'"] \\
Y & Y' \arrow[l, "p"] & ,
\end{tikzcd}
\]
with $f$ a $\D$-fibration, and $p$ a $\D$-equivalence, it does not necessarily follow that $q$ is a $\D$-equivalence. (But this will nevertheless be the case, ``of course'', if $\D = \HOT$.)

\msp{page 120}On the other hand, we have:

\textbf{Proposition 8.} \textit{Let $f : X \to Y$ be in \Dia. Equivalent conditions:}

a) \textit{$f$ is a $\D$-fibration.}

b) \textit{For every $p : Y'' \to Y'$ over $Y$, if $p$ is a} 1-\textit{complete $\D$-equivalence} (def. 1, p. 113), \textit{then $q : X'' \to X'$ is a $\D$-equivalence.}
\[
\begin{tikzcd}
X \arrow[d, "f"'] & X' \arrow[l] \arrow[d, "f'"'] & X'' \arrow[l, "q"'] \arrow[d, "f''"] \\
Y & Y' \arrow[l] & Y'' \arrow[l, "p"]
\end{tikzcd}
\]
\textbf{Corollary.} \textit{Suppose that in \Dia, every $\D$-equivalence is} 1-\textit{complete.} (N.B. I know only the derivator \HOT{} to satisfy this condition, which seems almost characteristic of \HOT.) \textit{Then $f : X \to Y$ is a $\D$-fibration if and only if for every $p : Y'' \to Y'$ over $Y$, if $p$ is a $\D$-equivalence, then $q : X'' \to X'$ too. In this exceptional case among derivators, the notion of $\D$-fibration is stable under composition} (which is tautological from the criterion I have just given.)

\printedpage{72}
% head: A. Grothendieck : Dérivateurs, Ch. I, Édition des Universités de Montpellier II et Paris VII

We are going to strengthen the notion of $\D$-fibration, to make it more flexible, and in particular stable under composition.

\textbf{Proposition 9.} \textit{Let $f : X \to Y$ be in \Dia. Equivalent conditions:}

a) \textit{For every $y \in Y$, $X_y \to X/y$ is a \underline{cohomological} $\D$-\underline{equivalence}} (def. 1, p. 113), \textit{and likewise for $X/y \to X/y'$ for every arrow $y \to y'$ in $Y$. And these conditions on $f$ remain valid after any base change} (N.B. it suffices that the first of these conditions remain so).

\msp{page 121}b) \textit{for every $F \in \Ob\ \D(X)$ \underline{locally} \underline{constant}, $f_*F$ is locally constant on $Y$ and is computed ``fiber by fiber'' -- and likewise after any base change.} % Doubt resolved: attached PDF leaf 36 (printed p.72) has an en dash, encoded --.

c) \textit{For every $p : Y'' \to Y'$ over $Y$, if $p$ is a cohomological $\D$-equivalence, then the same holds for $q : X'' \to X'$.}

c$'$) \textit{As in} c, \textit{but assuming as a hypothesis on $p$ that $Y''$ is the final category, and $Y'$ has a final object.}

\textsc{Proof}. In the order ($\pm$ tautological)
\[
\mathrm{a} \Longrightarrow \mathrm{b} \Longrightarrow \mathrm{c} \Longrightarrow \mathrm{c}' \Longrightarrow \mathrm{a} \quad .
\]

\textbf{Definition 5.} Under the equivalent conditions of proposition 9, we say that $f$ is a \textit{cohomological $\D$-fibration}. Under the dual conditions, we say that $f$ is a \textit{homological $\D$-fibration}. We say that $f$ is a \textit{strong $\D$-fibration}, if it is both a cohomological $\D$-fibration and a homological $\D$-fibration.

These notions strengthen the notion of $\D$-fibration. They have the further advantage of being \underline{stable} \underline{under} \underline{composition}, not only under base change.

\msp{page 122}The notion of strong $\D$-fibration was chosen to cover the two dual notions of cohomological and homological $\D$-fibration (for lack of a better notion) while remaining stable under composition. One could require that for every % [sic] meilleur -> meilleure, agreement with "notion"
\[
Y'' \longrightarrow Y'
\]
over $Y$, with $Y'$, $Y''$ \Hot-aspheric, $q : X'' \to X'$ be a 1-complete $\D$-equivalence (or even 2-complete), call these the 1-\textit{complete} (resp. 2-\textit{complete}) $\D$-\textit{fibrations}. This strengthens in a rather elegant way the notion of strong $\D$-fibration, but has the great fault of not being stable under composition (as far as I know!), except in the particular case where we have nice internal \underline{\Hom} (cf. bottom of page 113), in which case the notion of $\D$-equivalence already implies that of 1-complete $\D$-equivalence,
\printedpage{73}
% head: A. Grothendieck : Dérivateurs, Ch. I, Édition des Universités de Montpellier II et Paris VII
\msp{page 123}thus also that of homological $\D$-equivalence. So in this case
\[
\begin{tikzcd}
\text{cohomological $\D$-fibration} \arrow[r, Leftrightarrow] \arrow[d, Rightarrow] & \text{1-complete $\D$-fibration} \\
\text{homological $\D$-fibration} \quad , &
\end{tikzcd}
\]
and there is no need to look further than cohomological $\D$-fibrations.

One could think, it is true, of defining (in the general case) a stronger notion, by requiring that for every $p : Y'' \to Y'$ over $Y$ which is a 1-complete $\D$-equivalence, $q : X'' \to X'$ also be a 1-complete $\D$-equivalence. This time, the notion is stable under composition, but it nevertheless seems unusable to me (outside the case we have just mentioned, where it coincides with the notion of cohomological $\D$-fibration, or again of strong $\D$-fibration), for, not being in a position to describe the 1-complete $\D$-equivalences, I am also unable to give nontrivial examples of ``strictly 1-complete'' $\D$-fibrations.

To check that an $f : X \to Y$ in \Dia{} is a cohomological $\D$-fibration, one must
\msp{page 124}know how to check that arrows $q : Z' \to Z$ in \Dia{} are cohomological $\D$-equivalences, hence that $H^*_{\D}(Z,F) \overset{\sim}{\rightarrow} H^*_{\D}(Z',q^*(F))$ is \texttt{[an]} isomorphism for every $F \in \Ob\ \D(Z)$ which is locally constant on $Z$. Now I know how to check a statement of this nature (outside the case where I know in advance that $Z$ is completely $\D$-aspheric, \textit{i.e.} that the locally constant coefficients on $Z$ are in fact constant) \underline{only} if $q$ is already a \Hot-equivalence. % [sic] constant) -> constants), agreement with "coefficients"
That is to say that I know how to exhibit cohomological $\D$-fibrations, only when they are already cohomological \HOT-fibrations. I want, in what follows, to give usable criteria for \HOT-fibrations. But we are going to axiomatize these developments, keeping a general $\D$, on which I am going to make an assumption so draconian that practically, I know only the derivator \HOT{} which satisfies it. This is $c_1$ below.

\textbf{Der} 3 $c_1$) Every $\D$-equivalence is a cohomological $\D$-equivalence.

$c_2$) Every $\D$-equivalence is a homological $\D$-equivalence.

\msp{page 125}$c'$) Every $\D$-equivalence is 1-complete.

$c''$) Every $\D$-equivalence is 2-complete.

N.B. We have
\[
\begin{tikzcd}
& c'' \arrow[d, Rightarrow] & \\
& c' \arrow[dl, Rightarrow] \arrow[dr, Rightarrow] & \\
c_1 & & c_2
\end{tikzcd}
\]

\end{document}
